\chapter{Parent-origin портала радиуса стрелочной когерентности}
\label{ch:v8-baryon-c0-edge-coherence-radius-portal-parent-origin}

\section{Общий алгебраический носитель}

Предыдущая классификация выбрала
$T=\Tr(BB^*)$ как единственный активный составной singlet с уже
выведенным ненулевым значением. Спектральный parent когерентности использует
градуированную цепь
\begin{equation}
 \mathcal H_B=\mathcal H_0\oplus\mathcal H_1\oplus\mathcal H_2,
 \qquad
 \dim_{\mathbb C}(\mathcal H_0,\mathcal H_1,\mathcal H_2)=(1,6,3),
 \qquad
 \mathcal D_B=
 \begin{pmatrix}0&A_B^*&0\\A_B&0&C_B^*\\0&C_B&0\end{pmatrix}.
 \label{eq:v8-coherence-portal-chain}
\end{equation}
Совпадение чётных углов с размерностями $1+3$ позволяет сначала проверить
чисто алгебраическую гипотезу. Продолжение центрального
$Q=\operatorname{diag}(-3/4,(1/4)I_3)$ на всю цепь имеет вид
\begin{equation}
 \widehat Q(q)=\operatorname{diag}
 \left(-\frac34,qI_6,\frac14I_3\right).
 \label{eq:v8-coherence-portal-general-q-extension}
\end{equation}
Требование единого бесследового направления фиксирует середину:
\begin{equation}
 \Tr\widehat Q(q)=6q=0
 \quad\Longrightarrow\quad
 \widehat Q=\operatorname{diag}
 \left(-\frac34,0_6,\frac14I_3\right).
 \label{eq:v8-coherence-portal-unique-q-extension}
\end{equation}

Обозначая $T=\Tr(BB^*)$ и $d=\det(BB^*)$, точный compound-matrix расчёт
даёт требуемый смешанный след
\begin{equation}
 \Tr(\widehat Q\mathcal D_B^2)=-\frac58T.
 \label{eq:v8-coherence-portal-marked-second-moment}
\end{equation}

\section{Полные моменты общего оператора}

Положим
\begin{equation}
 X_{B,\lambda}=\mathcal D_B+\lambda\widehat Q.
 \label{eq:v8-coherence-portal-total-operator}
\end{equation}
Его первые моменты равны
\begin{equation}
 \Tr X_{B,\lambda}^2=3T+\frac34\lambda^2,
 \label{eq:v8-coherence-portal-second-moment}
\end{equation}
\begin{equation}
 \Tr X_{B,\lambda}^3
 =-\frac{15}{8}\lambda T-\frac38\lambda^3,
 \label{eq:v8-coherence-portal-cubic-moment}
\end{equation}
и
\begin{equation}
 \Tr X_{B,\lambda}^4
 =\frac94T^2+\frac{15}{4}d
 +\frac{19}{8}\lambda^2T+\frac{21}{64}\lambda^4.
 \label{eq:v8-coherence-portal-fourth-moment}
\end{equation}
Следовательно, любой parent, составленный только из унаследованных чётных
моментов, удовлетворяет
\begin{equation}
 \left.\partial_\lambda\Tr X_{B,\lambda}^{2}\right|_{\lambda=0}
 =
 \left.\partial_\lambda\Tr X_{B,\lambda}^{4}\right|_{\lambda=0}=0.
 \label{eq:v8-coherence-portal-even-moment-no-source}
\end{equation}
Он может создать член $\lambda^2T$, то есть изменить жёсткость щели на
фоне конденсата, но не создаёт линейный источник $\lambda T$.

\section{Условный кубический маршрут}

Если отдельно допустить коэффициент $c_3$ нечётного момента, то
\begin{equation}
 c_3\Tr X_{B,\lambda}^3
 =-\frac{15c_3}{8}\lambda T-\frac{3c_3}{8}\lambda^3.
 \label{eq:v8-coherence-portal-cubic-parent-decomposition}
\end{equation}
Таким образом, портал возникает с фиксированной связью с собственным
кубическим членом:
\begin{equation}
 \kappa_B=\frac{15}{8}c_3,
 \qquad
 c_{\lambda^3}=\frac38c_3,
 \qquad
 \frac{\kappa_B}{c_{\lambda^3}}=5,
 \qquad
 j_{\rm eff}\big|_{T=3}=\frac{45}{8}c_3.
 \label{eq:v8-coherence-portal-fixed-cubic-ratio}
\end{equation}
Real-завершение само по себе не уничтожает эту форму. Для сопряжённой пары
самосопряжённых операторов
\begin{equation}
 \frac12\left(\Tr X^3+\Tr\overline X^{,3}\right)=\Tr X^3\in\mathbb R.
 \label{eq:v8-coherence-portal-real-half-trace}
\end{equation}
Препятствие состоит не в Real-удвоении, а в отсутствии $c_3$ в текущем
чётном спектральном полиноме. Общий принцип спектрального действия задаёт
функционал от спектра полного оператора, но не разрешает добавлять нужный
нечётный момент после просмотра результата
\cite{v8-coherence-portal-spectral-action}.

\section{Типизационный запрет размерностного совпадения}

Есть и более ранний барьер. Трёхмерный угол coherence-chain равен
\begin{equation}
 \mathcal H_2=\Lambda^2V\otimes\Lambda^2W^*,
 \qquad
 W=W_{eX}\oplus W_Y\simeq\mathbb C^2\oplus\mathbb C,
 \qquad
 \mathcal H_2\simeq\mathbf1\oplus\mathbf2
 \label{eq:v8-coherence-portal-channel-corner-split}
\end{equation}
относительно фактической канальной группы
$U(2)_{eX}\times U(1)_Y$. Это не уже объявленный неприводимый семейный
триплет. Различие фиксируется коммутантами:
\begin{equation}
 \dim\operatorname{End}_{U(2)_{eX}\times U(1)_Y}(\mathcal H_2)=2,
 \qquad
 \dim\operatorname{End}_{SO(3)_{\rm fam}}(\mathbb C^3)=1.
 \label{eq:v8-coherence-portal-commutant-mismatch}
\end{equation}
Поэтому совпадение размерностей $1+3$ не задаёт канонического
типизированного отождествления двух носителей. Функториальная
superconnection-конструкция корректна на своём associated bundle
\cite{v8-coherence-portal-superconnection}, но не превращает канальные
миноры в семейные состояния.

Итоговые реестры равны
\begin{equation}
 N_{\rm algebraic\ shape}=4/4,
 \qquad
 N_{\rm parent\ origin}=0/2
 \quad\{\text{typed intertwiner},c_3\}.
 \label{eq:v8-coherence-portal-ledgers}
\end{equation}

\begin{theorem}[Алгебраическая форма портала из кубического момента]
Единственное бесследовое продолжение центрального направления на цепь
$1\to6\to3$ даёт
$\Tr(\widehat Q\mathcal D_B^2)=-5T/8$. Поэтому кубический момент общего
оператора содержит портал $-15\lambda T/8$ и собственный член
$-3\lambda^3/8$ с фиксированным отношением коэффициентов $5$.
\end{theorem}

\begin{nogocriterion}[Запрет размерностного и нечётного импорта]
Нельзя объявлять этот портал выведенным: coherence-угол имеет канальный тип
$1+2$, а не семейный тип $3$, и унаследованный чётный spectral parent не
содержит $c_3\Tr X^3$. Требуются отдельно типизированный intertwiner и
parent-origin нечётного коэффициента.
\end{nogocriterion}

\begin{protocol}\begin{enumerate}
\item общий алгебраический носитель размерности $10$: \textbf{построен};
\item бесследовое продолжение $Q$: \textbf{единственно};
\item смешанный след $\Tr(\widehat Q\mathcal D_B^2)$:
\textbf{$-5T/8$};
\item чётные моменты создают линейный портал: \textbf{нет};
\item кубический момент создаёт портал: \textbf{условно да};
\item отношение portal/self-cubic: \textbf{$5$};
\item Real-полуслед уничтожает форму: \textbf{нет};
\item coherence-угол является текущим family-triplet: \textbf{нет};
\item алгебраическая форма: \textbf{$4/4$};
\item parent-origin: \textbf{$0/2$};
\item следующий гейт: \textbf{intertwiner чётного угла и family-triplet}.
\end{enumerate}\end{protocol}

Машинный сертификат сохранён в
\path{s2t_v8_baryon_c0_singlet_triplet_central_gap_edge_coherence_radius_portal_parent_origin_gate_results.json}.

\begin{thebibliography}{9}
\bibitem{v8-coherence-portal-spectral-action}
A.~H.~Chamseddine, A.~Connes,
\emph{The Spectral Action Principle}, arXiv:hep-th/9606001.
\bibitem{v8-coherence-portal-superconnection}
G.~Roepstorff,
\emph{Superconnections and the Higgs Field}, arXiv:hep-th/9801040.
\bibitem{v8-coherence-portal-quiver-morse}
M.~Harada, G.~Wilkin,
\emph{Morse theory of the moment map for representations of quivers},
arXiv:0807.4734.
\end{thebibliography}